\scanpage{364}
\mspage{103}

\textcircled{b} $C_\infty$, et les objets de $\mathrm{Cat\,hot}$ et de $\mathrm{Hot}$ qui lui correspondent. Dans quelle mesure ces derniers sont ils indépendants du choix du syst. proj. cofinal à g. $(t_n)_{n\in\mathbb N}$ dans $T$ ? Quel est le rôle de $C_\infty$ ? Est-elle \inserted{canoniquement} isomorphe à $C$ dans $\mathrm{Cat\,hot}$, ou dans $\mathrm{Hot}$ ? (Auquel cas la question de son ``rôle'' me semblerait résolue !)

\textcircled{c} Le ind-objet $H_C$ dans $\mathrm{Cat\,hot}$, et \unsure{\deleted{le}|le} ind-objet de $\mathrm{Hot}$ qu'il définit.

Je note que le Ind-objet $H_C$, et aussi les catégories $C_\infty$ construites comme dessus, ne changent pas quand on remplace $T$ par une sous-catégorie \inserted{$T'$} cofinale à gauche, ou \deleted{par} aussi par une catégorie de la forme $T_{/t_0}$ \inserted{$(t_0\in\mathrm{Ob}\,T)$}, qui donne lieu à \deleted{\unsure{le} \illegible{1} \unsure{de|à} \unsure{cat.}} \inserted{l'ens. ordonné} \deleted{\unsure{ord}} associé
\[
\bigl(T_{/t}\bigr)_0=T_{0/t}=\{t^{\unsure{2|\tau|\prime|?}}\in T_0\mid \unsure{t_0|t}\leq t_0\}.
\]

Si la relation entre les \deleted{catégories} \inserted{objets} dans a), et ceux dans b) \unsure{ou} dans c), est très étroite, on en conclurait que le ``type'' \deleted{d'homotopie} de $C$ dans $\mathrm{Cat\,hot}$, ou dans $\mathrm{Hot}$ suivant les cas, ne change pas en remplaçant $C$ par

% archive sheet number: 363
